\documentclass[10pt,a4paper]{article}

\usepackage[margin=0.75in]{geometry}
\usepackage{amsmath}

% Tighten spacing to fit two pages.
\setlength{\parskip}{2pt}
\setlength{\parindent}{0pt}

% Minimal section spacing tweaks (no titlesec dependency).
\makeatletter
\renewcommand\section{\@startsection{section}{1}{\z@}{6pt}{2pt}{\normalsize\bfseries}}
\renewcommand\subsection{\@startsection{subsection}{2}{\z@}{4pt}{1pt}{\small\bfseries}}
\makeatother

% Compact title.
\makeatletter
\renewcommand\maketitle{
  \begin{center}
    {\large\bfseries\@title}\par\vspace{2pt}
    {\small\@author}\par\vspace{1pt}
    {\small\@date}\par
  \end{center}
  \vspace{-0.3em}
}
\makeatother

\title{MLSys 2026 Scheduling Contest -- Track A}
\author{Toni Boehnlein \quad P\'al Andr\'as Papp}
\date{April 24, 2026}

\begin{document}
\maketitle

\section{Overview}

We treat the scheduling problem as a two-level search: an outer
\emph{partition} of operations into subgraphs, and an inner \emph{coupling}
layer that stitches adjacent subgraphs together via fast-memory retention.
Both levels optimise the same roofline cost model, share a thread-safe
memoisation cache, and are driven by a Fiduccia--Mattheyses (FM) local search
embedded inside an evolutionary multi-start loop. Symmetry in the input DAG
(parallel branches, repeated series motifs) is detected once via Merkle
hashing and used to guide local search and crossover.

The binary targets the 2--120\,s timeouts in the contest harness: the time
budget is parsed from the benchmark file name and strictly honoured by a
single deadline propagated to every worker thread.

\section{Cost model}

The reference cost is re-implemented bottom-up in
\texttt{src/core/subgraph.cpp}. For a candidate op-set we derive the
tile-slice sizes by backward propagation from the boundary outputs, classify
each boundary tensor as \emph{once/row/col/tile/stream load} (based on
which tile-index component drives it), and then sum a per-tile roofline
$\max(\text{compute},\text{memory})$. Split-$k$ is handled exactly: the
first and last $k$-steps pay the full per-tile I/O and the output evict,
while the middle $(nk{-}2)$ steps only stream inputs. Snake traversal
(RowMajor / ColMajor) is implemented as closed-form loop counts, not a
per-tile simulation, so enumerating every feasible $(w,h,k,\text{snake})$
config is cheap.

Feasibility is a first-class concept. \texttt{Subgraph::create} refuses any
op-set that has mixed internal/external consumers of the same tensor,
inconsistent boundary-output dimensions, or a disconnected sub-DAG.
Tiling divisibility ($w \mid W_{\text{out}}$, etc.) is precomputed as a
fast-path intersection. These invariants are verified by 47
unit/integration tests covering each move and each cost-model branch.

\section{Partition search}

The partition layer minimises the sum of subgraph costs assuming no
retention. The state is a set of groups (op subsets); moves are
\textsc{merge}, \textsc{steal}, \textsc{recompute}, \textsc{eject},
\textsc{internal-eject}, and \textsc{split}, plus two tensor-level
operators restricted to the FM layer -- \textsc{tensor-merge} (fuse all
groups containing consumers of a shared input tensor) and
\textsc{tensor-extract} (pull just the consumer ops into a new group).
Each move validates acyclicity on the condensed group DAG before
evaluation (modified Kahn with OR-semantics for recomputed ops) and
ephemeral-gap preservation (mixed-consumer constraint).

\subsection{Greedy descent and FM}

A greedy descent and an FM pass share the same move generators and the
same active-set data structure; they differ in how they explore. Greedy
seeds the active set with every op, never locks, and only accepts
strictly improving moves -- monotone by construction and therefore
guaranteed to converge to a local minimum. We run it whenever we need
to pull a perturbed partition back into a valley: at the end of every
initialisation strategy (before FM) and as a refinement polish on the
best partition FM returns. An FM pass, by contrast,
forms an active set from a small random seed of ops and then iterates:
pop the best move across the active set, apply it, lock the moved op
so it cannot initiate again this pass, and expand the active set with
the neighbouring ops of the affected groups (which are re-evaluated
with the locked set in mind). We accept worsening moves up to a
\emph{drift budget} $\varepsilon \cdot C_0$ below the best snapshot
seen so far; once the cumulative gain drops beyond this bound the pass
aborts. The outer loop restarts each subsequent pass from the
running best partition, with a cosine-annealed drift fraction that
decays from wide exploration at pass~0 to tight refinement near the
end. Each pass also returns its perturbed end state so that the
evolutionary layer can reseed from a still-drifting partition rather
than always from a local minimum.

\subsection{Evolutionary pool}

Generation~0 seeds the pool from six structurally diverse initialisers
(trivial, chain+edge, seed+grow, rev-topo, tensor-aligned, random) plus a
symmetry-driven init that reuses detected parallel/series patterns. Each
seed is refined by greedy + FM. Subsequent generations pick mutation
($2/3$) or crossover ($1/3$) tasks. Mutations are structural (random
merges, splits, reassignments, block moves, tensor merges) and ignore
cost, providing the diversity that the gain-guided FM cannot.

Crossover uses agreement clustering \`a la Sanders--Schulz: ops co-grouped
in both parents form inherited clusters, the remainder is reassigned
greedily. We additionally implement \emph{DAG-cut} crossover which slices
along topologically compatible boundaries and splices in a donor
sub-partition, and a symmetry-aware variant that restricts donor
candidates to Merkle-equivalent regions.

Pool admission is diversity-aware: partitions are compared via adjusted
Rand index (ARI) on the contingency table, canonicalised by Merkle
hashes so structurally symmetric partitions are collapsed. Insertion
protects the best-cost entry, admits any partition that is either
strictly better or sufficiently far from the nearest neighbour.

\section{Symmetry detection}

ML graphs are full of exact repetition (parallel heads, MoE experts,
stacked transformer blocks) and reusing a good grouping across replicas
is cheaper than re-searching for it. We precompute a directed Merkle
fingerprint per op: a forward hash composed from the op's init hash
(type, base cost, I/O shapes) and the sorted forward hashes of its
predecessors, iterated to a fixed point, plus a symmetric backward
hash; the combined hash identifies ops whose full predecessor and
successor cones are isomorphic. Equivalence classes of the forward
hash are grown into parallel \texttt{SymmetricPattern}s via a two-phase
orbit-level merger with a top/bottom level constraint to prevent
cycles, and a separate scan of the topological order extracts
\texttt{SeriesPattern}s -- consecutive blocks of $K$ ops with matching
internal and inter-block wiring. These patterns feed four places in
the search: \texttt{symm\_init} replicates a solved partition from one
component to all the others; a coverage-weighted symmetric-mutation
operator edits every replica simultaneously so the child stays on the
symmetric manifold; DAG-cut crossover restricts donor slices to
Merkle-equivalent regions so the cut boundaries match; and
\texttt{merkle\_canonicalise} rank-matches ops within each class
before ARI distance, so partitions that differ only by a symmetry
permutation hash to distance~0.

\section{Retention and coupling}

Retention is modelled as a directed coupling edge $g_a \to g_b$ carrying a
set of retained tensors. Each group has at most one incoming and one
outgoing edge, so coupling chains are disjoint paths in the group DAG.
The cost of a coupled pair recomputes both endpoints with retention
context and caches the result in a second cost-cache tier
(\texttt{evaluate\_with\_context}).

The coupled-partition layer runs its own FM/evo loop with all eight
partition moves lifted to respect chains, plus five coupling-specific
moves: \textsc{couple}/\textsc{uncouple} add or remove a coupling edge
$g_a\to g_b$ carrying a tensor $t$; \textsc{retain-force-split} splits a
group at a bridge edge so that an internal tensor becomes a boundary
output and immediately couples the two halves; \textsc{force-retain}
splits the destination group so a coupling that would otherwise
overflow fast memory becomes feasible; and \textsc{ephemeral-fuse}
extracts a producer and one of its consumers into a fresh group where
the shared tensor is ephemeral, and couples that group to the remaining
consumer. Each move validates chain-level acyclicity via a lifted BFS
reachability check before any cost evaluation.

Phase~2 seeds the coupling pool from ordering heuristics (DFS
retain-score and beam-search ordering) applied to every partition in
the Phase~1 pool, plus a random-coupling variant for diversity;
Phase~3 lets the coupled FM evo explore from there until the deadline.
The final solution falls back to the uncoupled best partition if the
coupled candidate fails validation.

\section{Engineering notes}

\textbf{Cost cache.} Every partition move and every coupling move
triggers a subgraph cost evaluation, so the cache is the hot path of
the whole solver. We use a two-tier lock-free open-addressed table
(\texttt{src/core/cost\_cache.h}): the base tier keys $\text{hash}%
(\text{sorted op-set}) \to \texttt{CostResult}$ and serves all
Phase~1 evaluations; the retention tier extends the key with
(entering, retain) sets and serves the coupled search. Keys are stored
as 64-bit wyhash-style fingerprints with no backing op-index
vector -- a correctness shortcut justified by a per-lookup collision
probability below $10^{-9}$ at realistic populations -- so a lookup is
a single atomic load and writes are a CAS on a per-slot atomic hash.
All threads and all phases share one cache instance; Phase~2/3 rely on
Phase~1 having warmed the base tier. Hit rates are well above 90\,\%
after the first generation, and the retention tier falls through to
the base tier when entering and retain are both empty so coupled
evaluations automatically reuse Phase~1 work.

\textbf{Incremental indexing.} The partition maintains an op$\to$group
index updated in-place on every move and a compact \texttt{GroupDAG}
with $O(|E|)$ \texttt{merge\_creates\_cycle} checks. Dirty sets are
propagated from the affected groups to their DAG neighbours so only
impacted moves are re-evaluated. Every move is split into a pure
\texttt{eval\_*} that checks feasibility and cost without mutating
state and an \texttt{apply\_*} that commits; each FM pass additionally
takes a lightweight checkpoint of the group state so a worsening pass
can roll back to its best snapshot cleanly.

\textbf{Parallelism.} Init, FM, and evolution tasks run on
\texttt{hardware\_concurrency()} threads (8 on the contest host)
dispatched from a shared atomic task counter; task granularity
is a full FM outer loop, so lock contention on the diversity pool is
negligible. All shared state (cost cache, pool) is either lock-free or
coarse-grained mutex-guarded.

\textbf{Safety margins.} The budget is clamped to 95\,\% of the posted
timeout and propagated as an absolute deadline, so every FM pass and
cost evaluation checks wall time and exits cleanly before the harness
kills the process.

\section{Results}

On the five released benchmarks (\texttt{mlsys-2026-\{1,5,9,13,17\}})
the solver returns valid solutions within the allotted time. The
evolutionary loop keeps improving the best partition on the larger
instances up to the 60\,s deadline; on small instances the pool
converges quickly and the remaining budget is spent on coupling
refinement.

\end{document}
